\chapter{Page faults}
\label{CH:PGFAULTS}
%%

The RISC-V CPU raises a page-fault exception
when a virtual address is used that has no mapping
in the page table, or has a mapping whose \lstinline{PTE_V}
flag is clear, or a mapping whose permission bits
(\lstinline{PTE_R},
\lstinline{PTE_W},
\lstinline{PTE_X},
\lstinline{PTE_U})
forbid the operation being attempted.
RISC-V distinguishes three
kinds of page fault: load page faults (caused by load instructions),
store page faults (caused by store instructions),
and instruction
page faults (caused by fetches of instructions to be executed).  The
\lstinline{scause} register indicates the type of the
page fault and the \indexcode{stval} register contains the address
that couldn't be translated.

The combination of page tables and page faults is a powerful tool.
Page tables give the kernel a level of indirection between virtual and
physical addresses, so that the kernel can control the structure and
content of address spaces. Page faults allow the kernel to intercept
loads and stores and, by modifying the page table, specify on the fly
what data those references refer to. The kernel can use these
capabilities to increase efficiency: for example, copy-on-write fork
allows the kernel to transparently share memory between parent and
child, avoiding the cost of copying pages that neither write.
Application programmers can also benefit. One possibility is
memory-mapped files, where the kernel uses paging to cause a file's
content to appear in an application's address space, transparently
reading file blocks in response to page faults. Another is lazy memory
allocation, which allows a program to ask for a huge virtual
address space, but only to pay the cost of allocating physical memory
for the pages the program actually reads and writes. The xv6++ kernel
in this repository does not use page faults to implement lazy allocation;
unexpected user page faults cause the faulting process to be killed.

Before proceeding, please read \texttt{kernel/sysproc.cpp},
\texttt{sys\_sbrk}, \texttt{process::grow} in \texttt{kernel/process.cpp},
and the virtual memory allocator in \texttt{kernel/virtual\_memory.cpp}.
Also look at the user-trap handler in the interrupt manager, which handles
unexpected traps from user space
(see \texttt{kernel/interrupt\_manager.h}, \texttt{user\_interrupt}).

\section{Eager allocation}
\label{sec:eager}

When an application asks for memory by calling \lstinline{sbrk}, xv6++
allocates physical pages immediately and maps them into the process's
page table before returning to user space. This is \indextext{eager
allocation}: if allocation fails, \lstinline{sbrk} can return an error
to the application before the application starts using the new address
range.

Eager allocation keeps the kernel simple. Once \lstinline{sbrk}
returns successfully, ordinary loads and stores to the newly allocated
memory can proceed without further kernel involvement. The downside is
that the kernel may allocate many pages that the program never actually
uses. For example, if an application asks for a gigabyte of memory, the
kernel must allocate and zero 262,144 4096-byte physical pages before
returning from \lstinline{sbrk}.

\section{Code}

The system call {\tt sbrk(n)} grows (or shrinks if {\tt n}
is negative) a process's memory size by {\tt n} bytes, and returns
the start of the newly allocated region (i.e., the old size).
The kernel implementation is \lstinline{sys_sbrk}
(see \texttt{kernel/sysproc.cpp}, \texttt{sys\_sbrk}).

In xv6++, process memory growth is ultimately performed by the process
object (see \texttt{kernel/process.h}, \texttt{process}) via its
memory-growth routine (see \texttt{kernel/process.h}, \texttt{grow}).

For positive arguments to \lstinline{sbrk}, \lstinline{process::grow}
calls the virtual memory manager to allocate and map physical memory
immediately. In classic xv6, this work is done by \lstinline{growproc}
and \lstinline{uvmalloc}. In xv6++, the same responsibilities are split
between the process object and the virtual memory manager (see
\texttt{kernel/virtual\_memory.h}, \texttt{alloc}) and (see
\texttt{kernel/virtual\_memory.h}, \texttt{mappages}), which allocate
physical pages, zero them, and create user PTEs.

When a process loads or stores to a virtual address that
lacks a valid page-table mapping, the CPU will
raise a \indextext{page-fault exception}.
The user-trap handler detects this situation
(see \texttt{kernel/interrupt\_manager.h}, \texttt{user\_interrupt})
and treats it as an unexpected trap from user space. Current xv6++
does not allocate pages on demand in response to user page faults;
instead, the user-trap handler marks the process as killed.

If an application frees memory using {\tt sbrk}, \lstinline{sys_sbrk}
shrinks the process and the kernel unmaps the corresponding pages. In classic xv6
this is done by \lstinline{shrinkproc} and \lstinline{uvmdealloc}; in xv6++ the
virtual memory manager performs the actual unmapping and freeing work
(see \texttt{kernel/virtual\_memory.h}, \texttt{dealloc}) and (see
\texttt{kernel/virtual\_memory.h}, \texttt{unmap}).
The unmap code frees the physical memory referred to by each valid PTE.

Note that xv6++ uses a process's page table not just to tell the
hardware how to map user virtual addresses, but also as the only
record of which physical memory pages are allocated to that
process. That is the reason why freeing user memory requires
examination of the user page table.

\section{Real world: Copy-On-Write (COW) fork}

Many kernels (though not xv6++) use page faults to implement
\indextext{copy-on-write (COW) fork}. The {\tt fork} system call
promises that the
child sees memory whose initial content is the same as the parent's
memory at the time of the fork. One way to implement this is to copy
the entire memory of the parent to newly allocated physical memory for
the child; this is what xv6 does. Copying can be slow, and it
would be more efficient if the child could share the parent's physical
memory. A straightforward implementation of this would not work,
however, since it would cause the parent and child to disrupt each
other's execution with their writes to the shared stack and heap.

Copy-on-write fork causes parent and child to safely share physical
memory by appropriate use of page-table permissions and page faults.
The basic plan is for the parent and child to initially share all
physical pages, but for each to map them read-only (with the
\lstinline{PTE_W} flag clear). Parent and child can then read from the
shared physical memory. If either writes a shared page, the RISC-V CPU
raises a page-fault exception. A kernel supporting COW would respond
by allocating a new page of physical memory and copying the shared
page into that new page. Then kernel would change the relevant PTE in
the faulting process's page table to point to the copy and to allow
writes as well as reads, and then resume the faulting process at the
instruction that caused the fault. Because the PTE now allows writes,
the re-executed store instruction will execute without a fault, and
will modify a private copy of the page rather than the shared page.

Copy-on-write requires book-keeping
to help decide when physical pages can be freed, since each page can
be referenced by a varying number of page tables depending on the history of
forks, page faults, execs, and exits. This book-keeping allows
an important optimization: if a process incurs a store page
fault and the physical page is only referred to from that process's
page table, no copy is needed.

Copy-on-write makes \lstinline{fork} faster, since \lstinline{fork}
need not copy memory. Some of the memory will have to be copied
later, when written, but it's often the case that most of the
memory never has to be copied.
A common example is
\lstinline{fork} followed by \lstinline{exec}:
a few pages may be written after the \lstinline{fork},
but then the child's \lstinline{exec} releases
the bulk of the memory inherited from the parent.
Copy-on-write \lstinline{fork} eliminates the need to
ever copy this memory.
Furthermore, COW fork is transparent:
no modifications to applications are necessary for
them to benefit.

\section{Real world: Demand paging}

Yet another widely-used feature that exploits page faults is
\indextext{demand paging}.  In the \lstinline{exec} system call, xv6 loads all
of an application's text
and data into memory before starting
the application.  Since applications
can be large and reading from disk takes time, this startup cost can
be noticeable to users. To
decrease startup time, a modern kernel doesn't initially load
the executable file into memory, but just creates the user page table with
all PTEs marked invalid. The kernel starts the program running;
each time the program uses a page for the first time, a page
fault occurs, and in response
the kernel reads the content of the page from disk and
maps it into the user address space.  Like COW fork and lazy
allocation, the kernel can implement this feature transparently to
applications.

The programs running on a computer may need more memory than the
computer has RAM. To cope gracefully, the operating system may
implement \indextext{paging to disk}. The idea is to store only a
fraction of user pages in RAM, and to store the rest on disk in a
\indextext{paging area}. The kernel marks PTEs that correspond to
memory stored in the paging area (and thus not in RAM) as invalid. If
an application tries to use one of the pages that has been {\it paged
  out} to disk, the application will incur a page fault, and the page
must be {\it paged in}: the kernel trap handler will allocate a page
of physical RAM, read the page from disk into the RAM, and modify the
relevant PTE to point to the RAM.

What happens if a page needs to be paged in, but there is no free
physical RAM? In that case, the kernel must first free a physical page
by paging it out or {\it evicting} it to the paging area on disk, and
marking the PTEs referring to that physical page as invalid. Eviction
is expensive, so paging performs best if it's infrequent: if
applications use only a subset of their memory pages and the union of
the subsets fits in RAM. This property is often referred to as having
good locality of reference. As with many virtual memory techniques,
kernels usually implement paging to disk in a way that's transparent
to applications.

Computers often operate with little or no {\it free} physical memory,
regardless of how much RAM the hardware provides. For example, cloud
providers multiplex many customers on a single machine to use their
hardware cost-effectively. As another example, users run many
applications on smart phones in a small amount of physical memory. In
such settings allocating a page may require first evicting an existing
page. Thus, when free physical memory is scarce, allocation is
expensive.

Lazy allocation and demand paging are particularly advantageous when
free memory is scarce and programs actively use only a fraction of
their allocated memory. These techniques can also avoid the work
wasted when a page is allocated or loaded but either never used or
evicted before it can be used.

\section{Real world: Memory-mapped files}

Other features that combine paging and page-fault exceptions include
automatically extending stacks and \indextext{memory-mapped files},
which are files that a program maps into its address space using
the \texttt{mmap} system call so that the program can read and write
them using load and store instructions.

%%
\section{Exercises}
%%

\begin{enumerate}

\item Implement lazy allocation for user memory grown with
\lstinline{sbrk}. The kernel should increase the process size without
immediately allocating physical pages, then handle user page faults in
the lazily allocated range by allocating and zeroing a page and mapping
it into the faulting process's page table. Make sure deallocation,
process exit, fork, and ordinary user tests continue to work when some
pages in the address space have not yet been allocated. The provided
\texttt{lazytest} program requests a large address range and touches only
a few pages to check that allocation happens on demand.

\end{enumerate}
